\section{Problem Statement}
\label{sec:problem}

Let $U = [2^L]$ be a universe of $L$-bit keys.
Given a set $S \subseteq U$ of $n$ keys, a maximum query range length $\mathcal{L}$,
and a target false positive rate $\varepsilon$, the \emph{range emptiness problem} asks
to build a data structure that answers ``$[a, b] \cap S = \emptyset$?'' for intervals of
length $\leq \mathcal{L}$, with:

\begin{itemize}
  \item \textbf{No false negatives:} if $[a, b] \cap S \neq \emptyset$, the answer is always
    ``not empty''.
  \item \textbf{Bounded false positives:} if $[a, b] \cap S = \emptyset$, the answer is
    ``not empty'' with probability at most $\varepsilon$.
\end{itemize}

Let $Y = U \setminus S$ denote all keys not stored in the structure.

\subsection{Information-Theoretic Bounds}
\label{sec:bounds}

\citet{goswami2015approximate} establish the following bounds.

\begin{theorem}[Lower bound, {\citet[§2]{goswami2015approximate}}]
Any data structure solving the range emptiness problem requires at least
$n \bigl(\log_2(\mathcal{L} / \varepsilon) - c\bigr)$ bits
for an absolute constant $c > 0$.
\end{theorem}

\begin{remark}
The lower bound is worst-case over the input: it holds for every
set~$S \subseteq U$ with $|S| = n$, with no assumptions on the
structure or distribution of~$S$.
\end{remark}

\begin{remark}
If the input is restricted to a class
$\mathcal{C} \subsetneq \binom{U}{n}$, the combinatorial entropy
of~$\mathcal{C}$ may be below
$n \log_2(\mathcal{L}/\varepsilon)$, and a tighter lower bound
applies.
\end{remark}

\begin{definition}[Locality-preserving hash]
\label{def:lp-hash}
Let $r \geq \mathcal{L}$. A hash function $h\colon U \to [r]$ is
\emph{locality-preserving} if it partitions $U$ into blocks of size~$r$
and applies an independent cyclic shift to each block. For any interval
$[a, b]$ of length at most $\mathcal{L} \leq r$, the image $h([a, b])$
is the union of at most two intervals in~$[r]$
(see Figure~\ref{fig:soda-mechanism}).
\end{definition}

\begin{theorem}[Upper bound, {\citet[§3]{goswami2015approximate}}]
The lower bound of Theorem~1 is achieved up to an additive $O(n)$ term
via two layers:
\begin{enumerate}
  \item A locality-preserving hash (Definition~\ref{def:lp-hash})
    $h\colon U \to U'$ where $|U'| = r = n\mathcal{L}/\varepsilon$.

  \item An \textbf{Exact Range Emptiness (ERE)} structure over $S' = h(S) \subseteq U'$:
    a succinct structure with zero error and space $n \log_2(r/n) + O(n)$ bits.
\end{enumerate}
\end{theorem}

The resulting space is $n \log_2(\mathcal{L}/\varepsilon) + O(n)$ bits, or equivalently
$\log_2(\mathcal{L}/\varepsilon) + O(1)$ bits per key --- matching the lower bound.

The fingerprint length $K = \lceil\log_2(n\mathcal{L}/\varepsilon)\rceil$ controls
accuracy: increasing $K$ by one bit halves the false positive rate. After ERE compression
(which subtracts $\log_2 n$ for block indexing), the effective cost per key is
$K - \log_2 n = \log_2(\mathcal{L}/\varepsilon)$.

\subsection{The Role of the Hash Function}
\label{sec:hash-role}

The hash $h$ projects $U$ down to $U' = [r]$. Under this projection, the stored keys
$S$ map to fingerprints $S' = h(S)$, and the non-keys $Y = U \setminus S$ map to
$Y' = h(Y)$.

A false positive occurs when $Y'$ overlaps with $S'$: a query point $y \in Y$ has
$h(y) \in S'$, so the structure cannot distinguish it from a real key. We call the
pre-images of a stored fingerprint $x' \in S'$ its \textbf{phantom points} --- all keys
$y$ such that $h(y) = x'$. The fewer phantoms overlap with $Y'$, the lower the FPR.

The ideal hash would map $S'$ and $Y'$ to disjoint regions of $U'$ --- zero
overlap, zero false positives. In practice this is impossible within the space budget,
but the choice of $h$ determines how well $S'$ and $Y'$ separate for a given data
distribution. This motivates two directions explored in this work: alternative hash
function designs (Chapter~\ref{chap:hash}) and data-adaptive segmentation strategies
that bypass hashing entirely on dense regions (Chapter~\ref{chap:hybrid}).
